% Capstone 8: LLL. Port of «Мини-курсовая - LLL».
% The original names no sizes or times ("приемлемое время"); the guideline numbers in
% parentheses come from the generated requirements.tex (make bench), see bench/README.md.
\documentclass{ca-capstone}
\usepackage[nointegrals]{wasysym}   % \smiley of footnote 4
\capstonerefs{sagemath,hps}
\begin{document}
% Generated by `make bench CAPSTONE=lll` (taskgen.capstone report) from
% bench/results.json: reference time x slack on the reference machine.
% Do not edit: change bench/config.yaml and re-run the benchmark.
\newcommand{\ReqSlack}{10}
\newcommand{\ReqDim}{30}
\newcommand{\ReqBits}{128}
\newcommand{\ReqSeconds}{100}
\newcommand{\ReqKnapN}{48}
\newcommand{\ReqKnapSeconds}{5}

\capstonetitle{Решение задачи редукции базиса алгоритмом LLL}

\goal
Требуется реализовать базовую версию алгоритма Lenstra"--~Lenstra"--~Lovasz редукции базиса
алгебраической целочисленной решётки.

\functionality
\begin{enumerate}
  \item Метод Гаусса редукции базиса двумерной алгебраической целочисленной
    решётки\footnote{метод Гаусса даёт точное решение, так что редукцию базиса двумерных решёток
    требуется осуществлять с его помощью}
  \item Проверка \emph{LLL-приведённости} некоторого базиса\footnote{т.е.\ проверка свойств
    \emph{LLL-приведённого} базиса (условия на длину векторов и условия Ловаша)}
  \item Реализация алгоритма \emph{LLL}
  \item Применение разработанных реализаций для взлома учебных криптосистем с открытым ключом:
    ранцевой и конгруентной\footnote{см.\ п.п.~6.1--6.2 в~\cite{hps}}
\end{enumerate}

\requirements
Требования к реализации:
\begin{enumerate}
  \item корректность
  \item приемлемая скорость работы (см. комментарий ниже)
\end{enumerate}

\textbf{Корректность работы} можно проверять, прибегая к помощи сторонних решателей, например,
из~\cite{sagemath}.

В данной работе важна не скорость, а корректность реализации и соответствие требованиям. Однако
тестовые примеры должны выполняться за приемлемое время\footnote{хотя бы до конца сдачи работы
\smiley} (ориентир: \emph{LLL}-редукция базиса решётки размерности до~\ReqDim{} с элементами
до~\ReqBits~бит "--- не более \ReqSeconds~с, взлом ранцевой криптосистемы с $n = \ReqKnapN$
элементами "--- не более \ReqKnapSeconds~с).

\details
В своей реализации \textbf{можно пользоваться библиотеками} для работы с большими целыми
числами, а также пакетами линейной алгебры.

Рекомендуется использовать язык \emph{Python} и систему компьютерной алгебры
\emph{SageMath}~\cite{sagemath}.

Желательно сохранять промежуточные результаты решения, такие как координаты базисных векторов на
каждой итерации алгоритмов Гаусса и \emph{LLL}, артефакты проверки \emph{LLL-приведённости}
базиса, подлежащие решётки при взломе учебных криптосистем и проч.

Платформо-независимый, \textbf{переносимый код} писать \textbf{не требуется}, но во время сдачи
должна быть возможность проверить работу.

\schedule
Работа выполняется \textbf{индивидуально}. Срок выполнения "--- \textbf{месяц}.

Работа должна быть ваша: реализация должна быть \textbf{разработана и написана вами}, а не
собрана из кусков чужих работ.

Желательно хранить код в \emph{git}-репозитории на \emph{github}, \emph{bitbucket} или подобных
платформах. \textbf{Документация не требуется}.

\remarks
\begin{enumerate}
  \item Рекомендуется сначала реализовать метод Гаусса, затем проверку условий
    \emph{LLL-приведённости} и обобщить Гаусса до полноценной базовой версии \emph{LLL}
  \item Подлежащие решётки можно формировать и задавать вручную, но автоматизация этого процесса
    будет плюсом вашей работы
\end{enumerate}

\references
\end{document}
